\chapter{Fundamenta\c{c}\~ao Te\'orica e Trabalhos Relacionados}
\label{cap:referencial}

\section{Minera\c{c}\~ao de Dados Educacionais e Learning Analytics}
\label{sec:edm-la}

A \sigla{EDM}{Minera\c{c}\~ao de Dados Educacionais} dedica-se ao desenvolvimento de m\'etodos para explorar dados oriundos de contextos educacionais com o objetivo de compreender, monitorar e melhorar processos de aprendizagem \cite{romero2010, handbookedm}. De forma complementar, \sigla{LA}{Learning Analytics} enfatiza a medi\c{c}\~ao, coleta, an\'alise e relato de dados sobre aprendizes e seus contextos para apoiar decis\~oes pedag\'ogicas e institucionais \cite{siemens2012, baker2014}.

Neste trabalho, essas duas perspectivas aparecem de forma integrada. Do lado de EDM, o projeto constr\'oi atributos, testa modelos supervisionados e compara m\'etricas. Do lado de LA, o foco est\'a em converter resultados anal\'iticos em relat\'orios compreens\'iveis para uso escolar.

\section{Previs\~ao de desempenho escolar}
\label{sec:previsao-desempenho}

Prever desempenho escolar pode assumir diferentes formas: prever evas\~ao, reprova\c{c}\~ao, engajamento, necessidade de refor\c{c}o ou nota futura. Este trabalho escolhe como foco principal a previs\~ao da pr\'oxima nota, porque esse alvo preserva granularidade pedag\'ogica e permite interven\c{c}\~oes mais precoces. Em paralelo, a nota prevista \'e convertida em um alerta de risco, representado por nota futura abaixo de 6,0, o que aproxima a modelagem do uso institucional.

\section{Modelos de aprendizado de m\'aquina no projeto}
\label{sec:modelos}

O projeto avaliou modelos supervisionados adequados a dados tabulares, com hist\'orico irregular entre alunos, disciplinas e anos letivos. A escolha dos candidatos levou em conta quatro crit\'erios pr\'aticos: capacidade de capturar rela\c{c}\~oes n\~ao lineares, robustez diante de atributos heterog\^eneos, custo razo\'avel de treinamento e possibilidade de comparar o desempenho com regras simples de refer\^encia.

Algoritmos como regress\~ao linear e \'arvore de decis\~ao simples seriam mais f\'aceis de explicar, mas tenderiam a representar de forma limitada intera\c{c}\~oes entre nota anterior, tend\^encia, faltas, pagamentos e contexto escolar. M\'etodos como \sigla{SVM}{Support Vector Machine} poderiam ser usados, mas costumam exigir maior cuidado com escala, ajuste de hiperpar\^ametros e interpreta\c{c}\~ao em bases tabulares maiores. Redes neurais profundas, por sua vez, foram consideradas menos aderentes ao volume e \`a irregularidade das sequ\^encias dispon\'iveis. Modelos como CatBoost e LightGBM permanecem como alternativas futuras, mas n\~ao foram priorizados na vers\~ao final para manter a solu\c{c}\~ao mais simples, reproduz\'ivel e apoiada em bibliotecas j\'a consolidadas no ambiente do projeto.

\subsection{Random Forest}

O \textit{Random Forest}, proposto por Breiman \cite{breiman2001}, foi importante no projeto por servir como primeiro modelo tabular forte ap\'os a fase explorat\'oria. Sua principal vantagem est\'a na robustez em dados tabulares, na capacidade de capturar rela\c{c}\~oes n\~ao lineares e na boa toler\^ancia a atributos heterog\^eneos.

\subsection{LSTM}

As redes \sigla{LSTM}{Long Short-Term Memory}, introduzidas por Hochreiter e Schmidhuber \cite{hochreiter1997}, foram consideradas por causa da natureza temporal do problema. No entanto, a base real revelou sequ\^encias irregulares, entradas e sa\'idas de alunos e hist\'oricos de extens\~ao muito desigual. Esses fatores reduziram a ader\^encia dessa abordagem ao escopo institucional do trabalho.

\subsection{Gradient Boosting}

O \textit{Gradient Boosting}, formalizado por Friedman \cite{friedman2001}, mostrou-se especialmente adequado \`a solu\c{c}\~ao final por trabalhar bem com dados tabulares e atributos derivados. A varia\c{c}\~ao \textit{HistGradientBoosting} foi inclu\'ida por ser eficiente em bases tabulares e por lidar bem com padr\~oes n\~ao lineares. Na rodada final, ele foi o melhor candidato para o modo de alerta de risco, enquanto o \textit{Random Forest} permaneceu competitivo e foi selecionado para a regress\~ao no modo de previs\~ao de nota.

\subsection{Classifica\c{c}\~ao supervisionada e regras simples}

Uma decis\~ao importante foi manter a classifica\c{c}\~ao de risco como tarefa supervisionada, em vez de utilizar apenas regras fixas, como ``se a nota prevista for menor que 6, ent\~ao o aluno est\'a em risco''. Essa regra seria simples, mas perderia parte do comportamento temporal do problema. O aluno pode ter uma nota prevista numericamente pr\'oxima da m\'edia e, ainda assim, apresentar padr\~ao de queda, hist\'orico anterior fraco na disciplina, faltas acumuladas ou dist\^ancia negativa em rela\c{c}\~ao \`a turma. A classifica\c{c}\~ao aprende diretamente com exemplos hist\'oricos quais combina\c{c}\~oes desses sinais antecederam notas futuras abaixo de 6,0.

Isso n\~ao elimina as regras. No projeto, regras e pontua\c{c}\~oes continuam sendo usadas na camada de relat\'orios e prioriza\c{c}\~ao. A diferen\c{c}a \'e que o alerta principal de risco \'e estimado por um classificador treinado com dados reais, enquanto as regras ajudam a organizar a comunica\c{c}\~ao pedag\'ogica e a calibrar prioridades.

\section{Dados temporais, validade e m\'etricas}
\label{sec:temporalidade}

Um ponto cr\'itico em aplica\c{c}\~oes educacionais \'e a valida\c{c}\~ao temporal. Em vez de dividir aleatoriamente treino e teste, este trabalho adota separa\c{c}\~ao por anos letivos, evitando vazamento entre passado e futuro. Essa decis\~ao aproxima a avalia\c{c}\~ao do uso real da solu\c{c}\~ao, em que a escola precisa prever o pr\'oximo per\'iodo a partir do passado j\'a conhecido. A preocupa\c{c}\~ao com valida\c{c}\~ao temporal tamb\'em \'e discutida em estudos sobre avalia\c{c}\~ao de modelos em dados dependentes do tempo, nos quais uma divis\~ao aleat\'oria pode produzir estimativas otimistas de desempenho \cite{bergmeir2018}.

As m\'etricas escolhidas acompanham essa dupla natureza do problema. Para regress\~ao, o trabalho utiliza \sigla{MAE}{Mean Absolute Error}, que mede o erro m\'edio absoluto em pontos de nota, e \sigla{RMSE}{Root Mean Squared Error}, que penaliza erros maiores com mais intensidade. Para classifica\c{c}\~ao de risco, utiliza \textit{precision}, \textit{recall} e \sigla{F1}{F1-score}. A \textit{precision} indica, entre os alunos sinalizados, quantos realmente estavam em risco; o \textit{recall} indica, entre os alunos que realmente estavam em risco, quantos foram recuperados pelo modelo; e o F1 combina essas duas medidas. Essa combina\c{c}\~ao permite avaliar tanto proximidade num\'erica quanto capacidade de recuperar casos cr\'iticos.

\section{\'Etica, privacidade e interpretabilidade}
\label{sec:etica}

O uso de dados educacionais exige aten\c{c}\~ao especial \`a privacidade, \`a finalidade do tratamento e \`a interpreta\c{c}\~ao cuidadosa das sa\'idas. Discuss\~oes em LA e EDM mostram que prote\c{c}\~ao de dados e governan\c{c}a n\~ao s\~ao barreiras para an\'alise, mas condi\c{c}\~oes para seu uso institucionalmente leg\'itimo \cite{slade2013, gasevic2016, khalil2016}. No contexto brasileiro, a \sigla{LGPD}{Lei Geral de Prote\c{c}\~ao de Dados Pessoais} refor\c{c}a esse cuidado \cite{lgpd}.

Por isso, o projeto incorporou anonimiza\c{c}\~ao da base e assumiu, desde o in\'icio, que a sa\'ida do sistema n\~ao poderia ser um veredito definitivo sobre o aluno. Os relat\'orios foram desenhados para indicar fatores de alerta e recomenda\c{c}\~oes iniciais, preservando o papel humano na decis\~ao pedag\'ogica \cite{fiok2022}.

\section{Trabalhos relacionados}
\label{sec:relacionados}

Os trabalhos relacionados ao tema podem ser agrupados em tr\^es eixos. O primeiro re\'une revis\~oes gerais da \'area, como as de Romero e Ventura \cite{romero2010} e de Baker e Inventado \cite{baker2014}. O segundo eixo envolve estudos preditivos aplicados \`a educa\c{c}\~ao, sobretudo aqueles que combinam m\'ultiplas fontes de informa\c{c}\~ao. O terceiro eixo aborda interpretabilidade, governan\c{c}a e quest\~oes \'eticas.

Em rela\c{c}\~ao a esses trabalhos, este TCC se diferencia por articular toda a trajet\'oria de um projeto aplicado: tratamento do banco, refatora\c{c}\~ao arquitetural, valida\c{c}\~ao temporal, gera\c{c}\~ao de relat\'orios e ranking executivo. Em vez de comparar apenas algoritmos em um \textit{dataset} pronto, o trabalho documenta a constru\c{c}\~ao completa de uma solu\c{c}\~ao voltada ao contexto escolar.
